% Generated draft: requires mathematical and source review.
\documentclass[11pt]{article}
\usepackage[margin=25mm]{geometry}
\usepackage{fontspec}
\setmainfont{Latin Modern Roman}
\usepackage{amsmath,amssymb,mathtools,unicode-math}
\setmathfont{Latin Modern Math}
\usepackage{enumitem,needspace,xurl,hyperref}
\setlength{\parindent}{0pt}
\setlength{\parskip}{5pt}
\setlength{\emergencystretch}{3em}
\newcommand{\N}{\mathbb{N}}
\newcommand{\Z}{\mathbb{Z}}
\newcommand{\Q}{\mathbb{Q}}
\newcommand{\R}{\mathbb{R}}
\newcommand{\C}{\mathbb{C}}
\newcommand{\F}{\mathbb{F}}
\newcommand{\A}{\mathbb{A}}
\newcommand{\PP}{\mathbb{P}}
\newcommand{\E}{\mathbb{E}}
\newcommand{\eps}{\varepsilon}
\newcommand{\dd}{\,\mathrm{d}}
\begin{document}
\section*{UnsolvedMath ID: 30005042}
\textbf{Source problem code:} catalogue entry
\subsection{Definitions and mathematical statement}
Let \(I\) denote the ordered parameter set of the coupled branching process; its precise range is not specified in the supplied general statement. The offspring process \(X=(X(\lambda))_{\lambda\in I}\) takes values in \(\Z_{\geq 0}\) and has almost surely nondecreasing sample paths:
\[
\mathbb P\bigl(X(\lambda)\leq X(\mu)\text{ for all }\lambda,\mu\in I\text{ with }\lambda\leq\mu\bigr)=1.
\]
Write \(Z_n(\lambda)\) for the generation-\(n\) population in the coupled branching process mentioned in the source. The supplied excerpt does not specify its construction, initial population, or full hypotheses.

For orientation, under the usual Galton--Watson interpretation, at a parameter for which \(0<m(\lambda):=\E[X(\lambda)]<\infty\), the conventional normalized population is
\[
W_n(\lambda)=\frac{Z_n(\lambda)}{m(\lambda)^n}.
\]
This convention does not identify the normalization or convergence assertion of the missing source theorem. Denote the limiting process referred to in the source by \(W\).

\begin{enumerate}
\item Can the source's assumption \((H)\) be weakened to an \(X\log X\) condition while retaining the full conclusion of its Theorem~1? A possible pointwise meaning of such a moment condition is
\[
\E\!\left[X(\lambda)\log^+X(\lambda)\right]<\infty
\quad\text{for every }\lambda\in I,
\]
where \(\log^+x=\max\{\log x,0\}\) for \(x>0\), and \(0\log^+0=0\). This displayed condition is only a candidate interpretation: the supplied record does not specify whether pointwise, locally uniform, or increment conditions are intended. Neither \((H)\) nor Theorem~1 is included, so the hypotheses, parameter range, topology, and mode of the requested convergence cannot be recovered precisely.
\item Characterize the distribution of the limiting process \(W\). Can the distributional fixed-point equation mentioned in the source help characterize that distribution? The equation itself is absent from the supplied record; its terms, independence conditions, and any normalization or uniqueness requirements are consequently unresolved. The target is the law of the process, including its joint distributions across parameters.
\end{enumerate}

The supplied motivational example takes \(X\) to be a simple Poisson process on \([0,\infty)\), so that \(\E[X(\lambda)]=\lambda\). For \(V=\{1,\ldots,n\}\), attach independent uniform \([0,1]\) variables \(U_e\) to all possible edges and define
\[
E(G(n,p))=\{e:U_e<p\},\qquad 0\leq p\leq1.
\]
The source describes the coupled branching object as the local limit around vertex~\(1\) of \(G(n,\lambda/n)\), jointly in \(\lambda\geq0\), as \(n\to\infty\); at finite \(n\), this graph expression uses \(0\leq\lambda\leq n\). Its stated single-parameter limit has Poisson\((\lambda)\) offspring. The supplied excerpt does not specify the joint convergence framework. It also reports that \((H)\) holds in this example with \(\kappa=1\), citing independent Poisson increments, but the inequalities defining \((H)\) and the end of that discussion are missing. These are supplied contextual claims, not independently verified conclusions.
\subsection{Short English statement}
The source asks whether an \(X\log X\) assumption suffices for the stated convergence theorem for a coupled branching process, and whether a distributional fixed-point equation characterizes its limiting process. The supplied excerpt omits the original hypotheses, theorem, and equation, so a complete precise formulation remains unresolved. The supplied status is open; this is a statement-only draft requiring human mathematical review.
\subsection{Sources}
\begin{itemize}
\item Supplied local problem-registry record, permanent UnsolvedMath ID \(30005042\), source problem code \texttt{catalogue entry}, title: \emph{Weak Moment Conditions and Branching-Process Limit Laws}. The record attributes the problem to the Oberwolfach Reports collection; no individual proposer is supplied.
\item Supplied source citation: \emph{Population Dynamics and Statistical Physics in Synergy (2023)}, Oberwolfach Reports, DOI \texttt{10.4171/owr/2022/12}, pp.~577--656. \url{https://doi.org/10.4171/owr/2022/12}. The registry gives proposed year 2022. These bibliographic labels are retained as supplied.
\item Supplied literature reference: \emph{Parametrised branching processes}. \url{https://arxiv.org/abs/2106.01426}.
\item Literature assessment dated exactly \texttt{2026-08-22}: the supplied curation labels the problem open and states that weakening the moment and increment assumptions to an \(X\log X\) condition and characterizing the limiting process distribution remain open, and that general branching results do not settle the coupled-process question. This is supplied dated curation, not an independently verified literature assessment.
\item The registry labels the statement \texttt{corrected\_verified} and says it was reconstructed and checked against the official report. That is a supplied verification claim; the present draft does not independently verify it or the cited sources.
\end{itemize}
\end{document}
